% =====================================================================
%  ROUTER-RAG POSTER - A1 PORTRAIT
%  Compile with: pdflatex (works on Overleaf out of the box)
% =====================================================================

\documentclass[final]{beamer}

% Poster format
\usepackage[size=a1,orientation=portrait,scale=1.0]{beamerposter}

% Standard packages
\usepackage[utf8]{inputenc}
\usepackage[T1]{fontenc}
\usepackage[english]{babel}
\usepackage{lmodern}
\usepackage{microtype}
\usepackage{booktabs}
\usepackage{array}
\usepackage{graphicx}
\usepackage{tikz}
\usetikzlibrary{positioning, arrows.meta, shapes.geometric, calc, fit}
\usepackage{amsmath}

% =====================================================================
% COLORS
% =====================================================================
\definecolor{primarypurple}{RGB}{108, 52, 131}
\definecolor{lightpurple}{RGB}{232, 218, 239}
\definecolor{accentorange}{RGB}{202, 111, 30}
\definecolor{lightorange}{RGB}{250, 215, 160}
\definecolor{modelgray}{RGB}{170, 170, 170}
\definecolor{lightgray}{RGB}{220, 220, 220}
\definecolor{successgreen}{RGB}{34, 153, 84}
\definecolor{lightgreen}{RGB}{169, 223, 191}
\definecolor{tealblue}{RGB}{45, 128, 150}
\definecolor{lightteal}{RGB}{184, 226, 232}
\definecolor{darktext}{RGB}{40, 40, 40}

% =====================================================================
% BLOCK STYLING
% =====================================================================
\setbeamercolor{block title}{fg=white, bg=primarypurple}
\setbeamercolor{block body}{fg=darktext, bg=lightpurple!25}
\setbeamerfont{block title}{family=\sffamily, size=\large, series=\bfseries}
\setbeamerfont{block body}{family=\sffamily}

% Alerted blocks for take-aways
\setbeamercolor{block alerted title}{fg=white, bg=accentorange}
\setbeamercolor{block alerted body}{fg=darktext, bg=lightorange!25}
\setbeamerfont{block alerted title}{family=\sffamily, size=\large, series=\bfseries}

% No navigation
\setbeamertemplate{navigation symbols}{}

% Small helpers
\newcommand{\metric}[2]{\textbf{\color{primarypurple}#1} #2}
\newcommand{\good}[1]{\textbf{\color{successgreen}#1}}
\newcommand{\warn}[1]{\textbf{\color{accentorange}#1}}
\newcommand{\code}[1]{\texttt{#1}}

% =====================================================================
% TITLE
% =====================================================================
\title{Routing When to Retrieve}
\subtitle{A Spanish RAG Conversational Assistant with Learned Query Routing}
\author{Ignacio Gutierrez \textbullet\ Paula Guerrero}
\institute{Project presentation -- 19.05.2026}

\begin{document}
\begin{frame}[t]

% --- TITLE BAR ---
\begin{center}
  {\Huge\bfseries\sffamily Routing When to Retrieve}\\[0.4em]
  {\Large\sffamily A Spanish RAG Conversational Assistant with Learned Query Routing}\\[0.6em]
  {\large\sffamily Ignacio Gutierrez \textbullet\ Paula Guerrero \quad | \quad Project presentation, 19.05.2026}
\end{center}
\vspace{0.5em}
\hrule\vspace{0.8em}

% =====================================================================
% Motivation + Research Questions
% =====================================================================
\begin{block}{Motivation \& Research Questions}
\begin{columns}[T,onlytextwidth]

\begin{column}{0.48\textwidth}
  Retrieval-Augmented Generation (RAG) helps factual question answering, but
  it is not always useful. Conversational prompts such as greetings, jokes, or
  personal support do not need a corpus lookup. Always retrieving adds latency
  and can inject irrelevant evidence; never retrieving loses grounding on
  knowledge-heavy questions.

  \vspace{0.5em}
  This project builds a Spanish assistant that first decides whether a query
  needs retrieval, then answers either directly or with retrieved SQAC context.
\end{column}

\begin{column}{0.48\textwidth}
  \textbf{Research questions:}
  \begin{itemize}
    \item RQ1: Can a learned router separate factual SQAC questions from open conversational turns?
    \item RQ2: Does routing preserve RAG answer quality while reducing latency?
    \item RQ3: Which router training strategy works best with few examples?
    \item RQ4: Does self-feedback repair unsupported or irrelevant RAG answers?
  \end{itemize}
\end{column}

\end{columns}
\end{block}

% =====================================================================
% Pipeline
% =====================================================================
\begin{block}{System Pipeline}
  \centering
  \begin{tikzpicture}[
    every node/.style={font=\small\sffamily, align=center},
    databox/.style={rectangle, draw=primarypurple, fill=lightpurple,
      thick, rounded corners, minimum height=1.35cm, minimum width=3.25cm,
      inner sep=4pt, text=darktext},
    procbox/.style={rectangle, draw=accentorange, fill=lightorange,
      thick, rounded corners, minimum height=1.35cm, minimum width=3.25cm,
      inner sep=4pt, text=darktext},
    modelbox/.style={rectangle, draw=modelgray, fill=lightgray,
      thick, rounded corners, minimum height=1.25cm, minimum width=3.25cm,
      inner sep=4pt, text=darktext},
    metricbox/.style={rectangle, draw=successgreen, fill=lightgreen,
      thick, rounded corners, minimum height=1.25cm, minimum width=3.25cm,
      inner sep=4pt, text=darktext},
    arrow/.style={->, >=Stealth, ultra thick, black}
  ]
    \node[databox] (query) at (0,0) {User query\\ \footnotesize Spanish};
    \node[procbox] (router) at (5.2,0) {Query router\\ \footnotesize retrieve vs. skip};

    \node[metricbox] (retrieve) at (10.8,1.8) {Hybrid retrieval\\ \footnotesize BM25 + BGE-M3};
    \node[metricbox] (rerank) at (16.1,1.8) {Reranker\\ \footnotesize MiniLM cross-encoder};
    \node[databox] (context) at (21.4,1.8) {Top-3 SQAC\\ passages};

    \node[modelbox] (direct) at (13.5,-1.8) {Direct path\\ \footnotesize no corpus lookup};
    \node[modelbox] (mixtral) at (27,0) {Mixtral-8x7B\\ \footnotesize 4-bit instruct model};
    \node[procbox] (feedback) at (32.8,0) {Self-feedback\\ \footnotesize ISREL + ISSUP};
    \node[databox] (answer) at (38.3,0) {Final answer\\ \footnotesize Spanish};

    \draw[arrow] (query) -- (router);
    \draw[arrow] (router) -- node[above, font=\footnotesize\sffamily] {retrieve} (retrieve);
    \draw[arrow] (retrieve) -- (rerank);
    \draw[arrow] (rerank) -- (context);
    \draw[arrow] (context) -- (mixtral);
    \draw[arrow] (router) -- node[below, font=\footnotesize\sffamily] {skip} (direct);
    \draw[arrow] (direct) -- (mixtral);
    \draw[arrow] (mixtral) -- (feedback);
    \draw[arrow] (feedback) -- (answer);
    \draw[arrow, dashed, color=accentorange] (feedback.north) to[bend left=25]
      node[above, font=\footnotesize\sffamily, color=accentorange] {regenerate / re-retrieve once} (context.north);
  \end{tikzpicture}

  \vspace{0.4em}
  \footnotesize The router is the control point: it decides whether the expensive retrieval branch is worth running.
\end{block}

% =====================================================================
% Data + Router Training
% =====================================================================
\begin{block}{Data, Router Training \& Evaluation}
\begin{columns}[T,onlytextwidth]

\begin{column}{0.32\textwidth}
  \textbf{Datasets and labels}\\[0.3em]
  \small
  \begin{tabular}{lrr}
    \toprule
    \textbf{Split} & \textbf{SQAC} & \textbf{Chat} \\
    \midrule
    Router train & 15,036 & 15,036 \\
    Router val & 1,779 & 1,664 \\
    Final test & 85 & 85 \\
    \bottomrule
  \end{tabular}

  \vspace{0.4em}
  SQAC examples are labelled \code{retrieve}; ChatSubs / chitchat examples are
  labelled \code{skip}. The retrieval index contains \textbf{5,601} SQAC passages.
\end{column}

\begin{column}{0.32\textwidth}
  \textbf{Router approaches}\\[0.3em]
  \small
  \begin{itemize}\setlength\itemsep{0.18em}
    \item \textbf{Frozen LR:} multilingual MPNet embeddings + logistic regression.
    \item \textbf{SetFit:} contrastive few-shot tuning over the same task.
    \item \textbf{Fine-tune:} full encoder fine-tuning with Hugging Face Trainer.
    \item \textbf{Rule fallback:} lexical Spanish retrieval triggers and conversational patterns.
  \end{itemize}
\end{column}

\begin{column}{0.32\textwidth}
  \textbf{Metrics}\\[0.3em]
  \small
  \begin{itemize}\setlength\itemsep{0.18em}
    \item Routing: accuracy, macro F1, per-class F1.
    \item Answer quality: token F1 and multilingual BERTScore F1.
    \item Retrieval: Recall@3 and MRR when retrieval is triggered.
    \item Cost: mean latency and retrieval rate.
    \item Oracle policy: retrieve iff \code{always\_rag} beats \code{never\_rag} by BERTScore.
  \end{itemize}
\end{column}

\end{columns}
\end{block}

% =====================================================================
% Implementation Details
% =====================================================================
\begin{block}{Implementation Details}
\begin{columns}[T,onlytextwidth]

\begin{column}{0.48\textwidth}
  \textbf{Retrieval stack}\\[0.3em]
  \small
  \begin{itemize}\setlength\itemsep{0.18em}
    \item \textbf{Lexical signal:} BM25 over Spanish-stemmed, stopword-filtered SQAC passages.
    \item \textbf{Semantic signal:} normalized \code{BAAI/bge-m3} dense embeddings.
    \item \textbf{Hybrid score:} 0.4 BM25 + 0.6 dense cosine similarity.
    \item \textbf{Reranking:} \code{cross-encoder/ms-marco-MiniLM-L-6-v2}.
    \item \textbf{Context budget:} final generator receives the top 3 reranked passages.
  \end{itemize}
\end{column}

\begin{column}{0.48\textwidth}
  \textbf{Generation and self-feedback}\\[0.3em]
  \small
  \begin{itemize}\setlength\itemsep{0.18em}
    \item Generator: \code{mistralai/Mixtral-8x7B-Instruct-v0.1}, loaded with 4-bit quantization when CUDA is available.
    \item RAG prompt instructs the model to answer only from retrieved fragments.
    \item Critic prompt asks the same LLM for JSON checks: \code{isrel} (context relevance) and \code{issup} (answer support).
    \item Failing support triggers one grounded regeneration; irrelevant context can trigger one query rewrite and re-retrieval.
  \end{itemize}
\end{column}

\end{columns}
\end{block}

% =====================================================================
% Findings 1 and 2
% =====================================================================
\begin{columns}[t]

\begin{column}{0.48\textwidth}
  \begin{block}{Finding 1: Routing Preserves Quality at Lower Cost}
    \small
    \begin{tabular}{lrrrr}
      \toprule
      \textbf{Strategy} & \textbf{Token F1} & \textbf{BERTScore} & \textbf{Latency} & \textbf{Retrieve} \\
      \midrule
      Always RAG & 7.88 & 61.10 & 13.94s & 100\% \\
      Never RAG & 4.57 & 58.21 & 6.64s & 0\% \\
      Router SetFit & \textbf{8.12} & \textbf{61.22} & 10.05s & 48.8\% \\
      Router fine-tune & 7.64 & 61.09 & \textbf{9.98s} & 50.6\% \\
      \bottomrule
    \end{tabular}

    \vspace{0.5em}
    \textbf{\color{successgreen}Best quality: Router SetFit.}
    It slightly beats Always RAG on both token F1 and BERTScore while retrieving
    on only about half of the test set.

    \vspace{0.4em}
    \textbf{\color{accentorange}Main cost win: latency.}
    Router SetFit reduces mean latency from 13.94s to 10.05s, a \textbf{27.9\%}
    reduction relative to Always RAG.
  \end{block}
\end{column}

\begin{column}{0.48\textwidth}
  \begin{block}{Finding 2: The Router Learns the Label Task Quickly}
    \centering
    \begin{tikzpicture}[
      x=1cm,y=1cm,
      every node/.style={font=\footnotesize\sffamily, text=darktext}
    ]
      \draw[->, thick] (0,0) -- (11.6,0) node[right] {training size};
      \draw[->, thick] (0,0) -- (0,5.1) node[above] {macro F1};
      \foreach \y/\lab in {0/75,1/80,2/85,3/90,4/95,5/100} {
        \draw[lightgray] (0,\y) -- (11.1,\y);
        \node[left] at (0,\y) {\lab};
      }
      \foreach \x/\lab in {0.8/8,2.3/64,3.5/128,4.8/500,6.3/2500,7.8/5000,9.2/12500,10.5/15036} {
        \draw (\x,0) -- (\x,-0.08);
        \node[below, rotate=25] at (\x,-0.12) {\lab};
      }

      % Approximate learning-curve points from models/router/learning_curve.json
      \draw[ultra thick, color=primarypurple]
        (0.8,2.29) -- (2.3,3.69) -- (3.5,3.65) -- (4.8,4.01)
        -- (6.3,4.30) -- (7.8,4.48) -- (9.2,4.57) -- (10.5,4.59);
      \foreach \p in {(0.8,2.29),(2.3,3.69),(3.5,3.65),(4.8,4.01),(6.3,4.30),(7.8,4.48),(9.2,4.57),(10.5,4.59)}
        \fill[primarypurple] \p circle (0.09);

      \draw[ultra thick, color=successgreen]
        (0.8,4.53) -- (2.3,4.86) -- (3.5,4.86) -- (4.8,4.88)
        -- (6.3,4.89) -- (7.8,4.93) -- (9.2,4.92) -- (10.5,4.93);
      \foreach \p in {(0.8,4.53),(2.3,4.86),(3.5,4.86),(4.8,4.88),(6.3,4.89),(7.8,4.93),(9.2,4.92),(10.5,4.93)}
        \fill[successgreen] \p circle (0.09);

      \draw[ultra thick, color=accentorange]
        (0.8,0.22) -- (2.3,2.22) -- (3.5,4.66) -- (4.8,4.90)
        -- (6.3,4.93) -- (7.8,4.94) -- (9.2,4.95) -- (10.5,4.95);
      \foreach \p in {(0.8,0.22),(2.3,2.22),(3.5,4.66),(4.8,4.90),(6.3,4.93),(7.8,4.94),(9.2,4.95),(10.5,4.95)}
        \fill[accentorange] \p circle (0.09);

      \node[anchor=west, color=primarypurple] at (6.3,2.1) {Frozen LR};
      \node[anchor=west, color=successgreen] at (6.3,3.45) {SetFit};
      \node[anchor=west, color=accentorange] at (6.3,2.75) {Fine-tune};
    \end{tikzpicture}

    \vspace{0.4em}
    \small SetFit reaches \textbf{97.6\% macro F1 with only 8 examples per class}.
    Full fine-tuning catches up after 128 examples per class and reaches the
    highest validation score at 12,500 examples per class.
  \end{block}
\end{column}

\end{columns}

% =====================================================================
% Findings 3 and 4
% =====================================================================
\begin{columns}[t]

\begin{column}{0.48\textwidth}
  \begin{block}{Finding 3: Label Accuracy Is Not the Same as Oracle Utility}
    \small
    \begin{tabular}{lrrr}
      \toprule
      \textbf{Strategy} & \textbf{Decision F1} & \textbf{Oracle acc.} & \textbf{Retrieve} \\
      \midrule
      Always RAG & 33.33 & 69.41 & 100.0\% \\
      Never RAG & 33.33 & 30.59 & 0.0\% \\
      Frozen LR router & 97.06 & 67.06 & 48.2\% \\
      SetFit router & 98.82 & 69.41 & 48.8\% \\
      Fine-tuned router & \textbf{99.41} & 69.41 & 50.6\% \\
      \bottomrule
    \end{tabular}

    \vspace{0.5em}
    The classifier almost perfectly recovers the dataset labels. The harder
    question is whether retrieval actually improves the generated answer for
    each SQAC sample.

    \vspace{0.4em}
    \textbf{\color{accentorange}Oracle ceiling is noisy.}
    On the SQAC subset, Always RAG beats Never RAG on only 69.4\% of samples by
    BERTScore. The best routers match that oracle-policy accuracy while using
    roughly half as much retrieval.
  \end{block}
\end{column}

\begin{column}{0.48\textwidth}
  \begin{block}{Finding 4: Self-Feedback Is Active but Still Conservative}
    \small
    \begin{tabular}{lrrrr}
      \toprule
      \textbf{Strategy} & \textbf{Trigger} & \textbf{Improved} & \textbf{Worsened} & \textbf{$\Delta$ Token F1} \\
      \midrule
      Always RAG & 81.2\% & 42.9\% & 25.9\% & +1.23 \\
      Router SetFit & 32.9\% & 11.8\% & 11.8\% & +0.08 \\
      Router fine-tune & 34.1\% & 12.4\% & 11.2\% & +0.02 \\
      Oracle SetFit 64 & 34.7\% & 16.5\% & 11.8\% & +0.60 \\
      \bottomrule
    \end{tabular}

    \vspace{0.5em}
    \textbf{\color{successgreen}Self-feedback helps most when retrieval is frequent.}
    Always RAG gives the critic more chances to repair unsupported answers.

    \vspace{0.4em}
    \textbf{\color{primarypurple}Routers reduce both errors and repair opportunities.}
    The best router skips 87 of 170 samples; among retrieved samples, the
    common actions are pass, regenerate, and direct fallback.
  \end{block}
\end{column}

\end{columns}

% =====================================================================
% Qualitative Analysis
% =====================================================================
\begin{block}{Qualitative Analysis --- What the Router Changes}
\begin{columns}[T,onlytextwidth]

\begin{column}{0.32\textwidth}
  \fcolorbox{successgreen}{lightgreen!30}{\parbox{0.95\linewidth}{
    \small
    \textbf{\color{successgreen}Good skip: conversational query}\\
    \textit{Query:} ``Dime si puedes estornudar''\\
    \textit{Route:} skip retrieval\\
    \textit{Effect:} the assistant answers directly instead of forcing an
    unrelated SQAC passage into a personal / conversational turn.
  }}
\end{column}

\begin{column}{0.32\textwidth}
  \fcolorbox{primarypurple}{lightpurple!30}{\parbox{0.95\linewidth}{
    \small
    \textbf{\color{primarypurple}Good retrieve: factual QA}\\
    \textit{Query:} ``¿Qué cargos ocupa Jesús Gil?''\\
    \textit{Route:} retrieve SQAC passages\\
    \textit{Answer:} ``alcalde de Marbella y presidente del Atlético de Madrid''\\
    \textit{Effect:} retrieved context supports a concise gold-like answer.
  }}
\end{column}

\begin{column}{0.32\textwidth}
  \fcolorbox{accentorange}{lightorange!30}{\parbox{0.95\linewidth}{
    \small
    \textbf{\color{accentorange}Remaining failure: bad evidence}\\
    \textit{Query:} ``¿Cuántos prisioneros han sido trasladados al hospital?''\\
    \textit{Gold:} ``doce''\\
    \textit{Route:} retrieve, then regenerate\\
    \textit{Effect:} the critic detects weak support, but retrieval did not
    recover the gold passage.
  }}
\end{column}

\end{columns}

\vspace{0.6em}
{\setlength{\fboxrule}{2pt}%
\fcolorbox{successgreen}{lightgreen!25}{\parbox{0.98\linewidth}{
  \vspace{0.2em}
  \textbf{\color{successgreen}Interpretation:} the router solves the broad
  interaction problem: do not retrieve for casual dialogue; retrieve for
  corpus-grounded QA. The remaining bottleneck is more often retrieval quality
  and evidence support than the binary routing decision itself.
  \vspace{0.2em}
}}}
\end{block}

% =====================================================================
% Take-aways
% =====================================================================
\vspace{0.5em}

\begin{alertblock}{Take-aways}
\textbf{\color{successgreen}A learned router is worthwhile.}
Router SetFit achieves the best measured answer quality while cutting retrieval
from 100\% to 48.8\% and reducing latency by 27.9\%.

\vspace{0.45em}
\textbf{\color{primarypurple}Few-shot SetFit is the practical sweet spot.}
It performs strongly with very little labelled data and is competitive with
full fine-tuning in the final pipeline.

\vspace{0.45em}
\textbf{\color{accentorange}Routing accuracy alone overstates success.}
The router reaches 98--99\% decision F1, but end-to-end answer gains depend on
retrieval quality, generator behavior, and whether retrieval actually helps a
specific sample.

\vspace{0.45em}
\textbf{\color{successgreen}Self-feedback is useful but not magic.}
It can repair unsupported answers, yet it cannot fix missing evidence when the
retriever fails to recover the right passage.
\end{alertblock}

\vspace{0.7em}
\hrule
\vspace{0.5em}
{\small\itshape \textbf{Limitations:} final test set is small (85 SQAC + 85 chat samples); Recall@3 and MRR are 0 in these runs, suggesting an index / ID alignment or retrieval-recall issue that deserves follow-up; oracle policy labels use BERTScore differences rather than human judgments; all generation results depend on one Mixtral prompt setup.}

% =====================================================================
% REFERENCES
% =====================================================================
\vspace{0.5em}
\hrule\vspace{0.4em}
\footnotesize\textbf{References:}\quad
Lewis et al.\ (2020).\ Retrieval-Augmented Generation.\ \emph{NeurIPS} \textbullet\
Reimers \& Gurevych (2019).\ Sentence-BERT.\ \emph{EMNLP} \textbullet\
Zhang et al.\ (2020).\ BERTScore.\ \emph{ICLR} \textbullet\
Izacard et al.\ (2021).\ Leveraging Passage Retrieval.\ \emph{EACL} \textbullet\
Mistral AI (2024).\ Mixtral of Experts.\ \emph{arXiv} \textbullet\
Karpukhin et al.\ (2020).\ Dense Passage Retrieval.\ \emph{EMNLP} \textbullet\
Schick \& Schütze (2021).\ Generating Datasets with Few-shot Prompts.\ \emph{EMNLP}

\end{frame}
\end{document}
